\chapter{Anti-Streptolysin O (ASO) Titer}

\section*{What Is This Test?}
The anti-streptolysin O test measures antibodies produced after infection with group A Streptococcus. It helps document a recent streptococcal infection when evaluating delayed complications rather than diagnosing an acute sore throat.

\section*{Alternative Names}
ASO; ASOT; Anti-Streptolysin O Antibody; Streptococcal Serology.

\section*{Why This Test Is Ordered}
Testing may support evaluation of suspected acute rheumatic fever or post-streptococcal glomerulonephritis when the preceding throat or skin infection is no longer detectable. It is not a substitute for a rapid strep test or throat culture in acute pharyngitis.

\section*{How This Test Is Performed}
A blood sample is tested for antibodies against streptolysin O. Paired samples collected several weeks apart may be more informative than a single result because antibody levels rise after infection and then decline.

\section*{Preparation, Risks, and Limitations}
No fasting is required. Venipuncture is the usual risk. ASO may remain elevated for months, may be lower after skin infection than after throat infection, and can be absent despite a clinically important streptococcal infection.

\section*{Understanding Results}
An elevated or rising titer supports recent streptococcal exposure but does not prove that current symptoms are caused by active infection. Interpretation depends on age, local reference ranges, clinical criteria, and complementary tests.

\section*{References}
World Health Organization guidance on rheumatic fever and post-streptococcal disease.

\clearpage
\section*{Understanding Results and Follow-up}
The laboratory report should identify the specimen, method, units, reference interval or decision limit, and whether the result is positive, negative, abnormal, or inconclusive. Reference intervals vary with age, sex, pregnancy, specimen type, assay, and laboratory; a result outside a range is not automatically a diagnosis, and a result within a range does not always exclude disease. Discuss unexpected results with the ordering clinician, who can decide whether repeat or confirmatory testing, treatment, or referral is appropriate. Do not change medicines or delay urgent care based on an isolated result.


\section*{Regulatory Status and Authorization Date}
Regulatory status applies to the specific assay, instrument, manufacturer, and intended use rather than to the test name alone. No single approval date can be assigned to this test category. For a commercial assay, verify the current product in the relevant regulator's database: the U.S. Food and Drug Administration (FDA) clearance, approval, or authorization; India's Central Drugs Standard Control Organisation (CDSCO) licence or permission; the European Union In Vitro Diagnostic Regulation (EU IVDR) CE certificate issued through the applicable conformity-assessment route; or the United Kingdom Medicines and Healthcare products Regulatory Agency (MHRA) registration and UKCA/CE status. Laboratory-developed tests may not have an individual FDA, CDSCO, EU IVDR, or MHRA approval date.
\clearpage
